\chapter{Functions and the stack}
\MARSChapterHeading{Functions and the stack}
\label{revision:functions-and-stack}

A loop gives several iterations one body of instructions. A function lets several
parts of a program use one body and then continue at their own call sites. That
continuation is part of the problem. A branch can reach the shared instructions,
but the shared instructions also need to know where to return.

We will extract a small operation from the sprite project: reverse the eight bits
of a row byte. Then we will call that operation twice to reverse two rows. The
second program must keep the first result while computing the second, and it must
return to its own caller after making those calls. These requirements explain why
some values can stay in temporary registers and others need a saved location.

This chapter assumes byte masks, loads, stores and compact branches from the
preceding chapters. The reversal introduces logical shifts and bitwise OR before
we use them to follow a call. By its end, you should be able to state a function's
inputs and effects, trace a call and return, distinguish temporary values from
values that must survive another call, and explain a stack frame byte by byte.

\section{Give a function a contract}

A function's label names its entry instruction. The label does not declare
parameters, allocate local variables or arrange a return. We need a convention
for those jobs. For these examples, arguments arrive in \texttt{\$a0} through
\texttt{\$a3}, and results leave in \texttt{\$v0} and, when needed,
\texttt{\$v1}. The names help us follow that agreement; they do not prevent an
instruction from writing a different register.

Our first function, \texttt{reverse\_byte}, uses the low eight bits of
\texttt{\$a0}. It returns their reversal as an integer from zero through 255 in
\texttt{\$v0}. Upper input bits do not affect the result. It performs no console
or device I/O and writes no memory. For example, input \texttt{0x16} has low-byte
bits \texttt{00010110}; the returned bits are \texttt{01101000}, or
\texttt{0x68}, decimal 104. Input \texttt{0x81} is unchanged by reversal because
its first and last bits are both one and its other bits are zero.

In the later sprite, bit seven supplies the leftmost pixel. Reversal exchanges
the two ends, then the next pair, and so on. We can test this byte operation
before drawing the image.

The function also promises to preserve the saved registers and stack pointer in
the convention below. A result alone is not the whole contract. A routine that
returns 104 while losing its caller's output address is unusable by that caller.

\section{A call records a continuation}

The instruction \texttt{balc reverse\_byte} transfers to the function and writes
a return address into \texttt{\$ra}, register 31. For \texttt{balc}, this is the
address of the machine instruction immediately after the call. If the call's
address is $C$, the saved address is $C+4$. The function returns with
\texttt{jrc \$ra}, the MARS assembler alias for \texttt{jic \$ra, 0}.
That instruction transfers to the address held in the register.

These are Release 6 compact transfers: they have no delay slot. The instruction
after \texttt{balc} waits until the function returns; it does not execute before
the function starts. The instruction after \texttt{jrc} is not a return delay
slot either. This differs from the retained delayed call/return forms you may
find in older MIPS examples.\footnote{The \emph{MIPS64 Instruction Set Reference
Manual}, revision 6.05, describes \texttt{balc} and its next-instruction link on
page 66. The worked sources use MARS's Release 6 compact forms throughout.}

The call transfers control and supplies a link. It does not copy an argument,
reserve stack space, or save every register. Preparing \texttt{\$a0} remains the
caller's job. Computing \texttt{\$v0} and reaching the return instruction remain
the function's jobs. When the same function is called at another address, the new
call records a different continuation without changing the function's body.

\section{Reverse the byte in a leaf function}

A function that calls no other function is a \emph{leaf function}. This one needs
no memory for local state: its remaining input, result and count fit in temporary
registers. The listing is a function body. The complete program later in the
chapter supplies its caller and entry point.

The shift instructions move bits between positions. In
\texttt{sll \$v0, \$v0, 1}, the destination comes first, the source second,
and the final operand says to shift left by one position. A zero enters at the
right; a bit shifted beyond the 32-bit word is discarded.
\texttt{srl \$t0, \$t0, 1} shifts the word right by one, discards its low bit
and fills the vacated high position with zero. These are word operations, whose
32-bit result is sign extended to register width. Our masked input and partial
reversal stay between zero and 255, so the upper register bits remain zero.

The \texttt{or} instruction combines two registers bit by bit: an output bit is
one if either corresponding input bit is one. Here the left shift has made room
by clearing the result's low bit. The \texttt{andi} selects the remaining input's
low bit into \texttt{\$t2}, and \texttt{or \$v0, \$v0, \$t2} puts that bit
in the cleared position without changing the other result bits.

\noindent\begin{minipage}{\linewidth}
\MARSCodeLabel{A leaf routine with one register argument and result}
\begin{minted}{mars-idea}
reverse_byte:
    andi $t0, $a0, 0xff
    move $v0, $zero
    li $t1, 8
reverse_bit:
    sll $v0, $v0, 1
    andi $t2, $t0, 1
    or $v0, $v0, $t2
    srl $t0, $t0, 1
    addiu $t1, $t1, -1
    bnezc $t1, reverse_bit
    nop
reverse_byte_return:
    jrc $ra
\end{minted}
\end{minipage}

The initial mask selects exactly the low byte, including when the original
register has nonzero upper bits. Each iteration makes room at the right of the
result by shifting it left, copies the remaining input's low bit into that room,
and shifts the remaining input right. The count makes this an eight-bit reversal,
not a loop that stops as soon as the remaining input becomes zero. Leading zeros
in a fixed-width pattern still need positions in the reversed result.

For \texttt{0x16}, the first four completed iterations have these states. All
values in the last two columns are hexadecimal.

\begin{center}
\begin{tabular}{rlll}
\toprule
Bits consumed & Bit just copied & Remaining input & Result \\
\midrule
0 & --- & \texttt{16} & \texttt{00} \\
1 & 0 & \texttt{0B} & \texttt{00} \\
2 & 1 & \texttt{05} & \texttt{01} \\
3 & 1 & \texttt{02} & \texttt{03} \\
4 & 0 & \texttt{01} & \texttt{06} \\
8 & 0 & \texttt{00} & \texttt{68} \\
\bottomrule
\end{tabular}
\end{center}

After $k$ iterations, the result contains the reversal of the $k$ bits already
consumed, in its low $k$ positions. The remaining input is the original low byte
shifted right by $k$, and the count is $8-k$. At eight iterations, all eight
positions have contributed. The result is at most 255, so these word shifts and
logical operations do not encounter a signed word boundary.

The \texttt{nop} after \texttt{bnezc} supplies a safe forbidden-slot instruction
on the final, not-taken branch. On taken iterations it is skipped. Once the loop
finishes, \texttt{jrc \$ra} returns with the result already in \texttt{\$v0}.
Because this routine never makes another call, no instruction in it replaces the
incoming return address.

\section{Decide which values must survive a call}

Now give \texttt{reverse\_pair} two inputs: the low byte of \texttt{\$a0} and the
low byte of \texttt{\$a1}. It returns their reversals in \texttt{\$v0} and
\texttt{\$v1}, respectively. For inputs \texttt{16} and \texttt{81} in
hexadecimal, those results must be \texttt{68} and \texttt{81}, in that order.

Before its first call, the pair routine needs to retain the second input. Before
its second call, it needs to retain the first result. Keeping that result only in
\texttt{\$v0} would fail: the second call is meant to replace \texttt{\$v0} with
its own result. Keeping the second input in \texttt{\$a1} happens to work with
today's leaf implementation, but its contract does not promise to retain argument
registers. Relying on an unpromised detail would tie the caller to that particular
implementation.

We use this small convention for the book's routines:

\begin{center}
\begin{tabular}{lp{0.69\linewidth}}
\toprule
Register group & Agreement between caller and callee \\
\midrule
\texttt{\$a0}--\texttt{\$a3} & Arguments; their values need not survive a call. \\
\texttt{\$v0}, \texttt{\$v1} & Results; another call can replace them. \\
\texttt{\$t0}--\texttt{\$t9} & Temporary work; their values need not survive a call. \\
\texttt{\$s0}--\texttt{\$s7}, \texttt{\$fp}, \texttt{\$gp} & A callee must return them with their incoming values. \\
\texttt{\$sp} & A callee must restore its incoming stack pointer. Call boundaries use 16-byte alignment here. \\
\texttt{\$ra} & A call replaces it. Retain an incoming return address if another call would destroy the continuation you still need. \\
\bottomrule
\end{tabular}
\end{center}

A \emph{caller} makes a call; its \emph{callee} receives that call. The same
routine can have both roles. \texttt{reverse\_pair} is called by \texttt{main},
then becomes a caller of \texttt{reverse\_byte}. It can keep values in saved
registers across those inner calls, but must first preserve the values its own
caller already had there.

This is an explicit convention for these standalone MARS programs, not a complete
operating-system ABI. An ABI also specifies matters such as register assignment,
stack arguments and linkage between separately built components. Follow that
component's ABI when interoperating with it; do not infer one from these register
names alone.

\section{Give the saved values a stack frame}

The stack is memory managed through \texttt{\$sp}. For this program it grows toward
lower addresses. A function reserves a \emph{stack frame} by subtracting its size
from \texttt{\$sp}, then uses offsets from the new pointer for its own slots.
Before returning, it restores the saved values and adds back exactly that size.
An inner call can reserve a lower frame without overlapping the outer one.

We need to save the incoming \texttt{\$ra}, \texttt{\$s0} and \texttt{\$s1}.
Each is a 64-bit value, so each slot occupies eight bytes. Three slots need 24
bytes; we reserve 32 to keep a 16-byte-aligned pointer aligned. The unused first
eight bytes are padding, not another saved register.

\begin{center}
\begin{tabular}{ll}
\toprule
Byte offsets from the new \texttt{\$sp} & Contents while this frame is active \\
\midrule
0--7 & Unused padding \\
8--15 & Incoming \texttt{\$s1} \\
16--23 & Incoming \texttt{\$s0} \\
24--31 & Incoming \texttt{\$ra} \\
32 and above & Outside this frame; belongs to the caller \\
\bottomrule
\end{tabular}
\end{center}

\MARSCodeLabel{A non-leaf routine preserves its caller's state}
\begin{minted}{mars-idea}
reverse_pair:
    daddiu $sp, $sp, -32
    sd $ra, 24($sp)
    sd $s0, 16($sp)
    sd $s1, 8($sp)
pair_frame_ready:
    move $s0, $a1
    balc reverse_byte
pair_first_ready:
    move $s1, $v0
    move $a0, $s0
pair_second_call:
    balc reverse_byte
pair_second_ready:
    move $v1, $v0
    move $v0, $s1
    ld $s1, 8($sp)
    ld $s0, 16($sp)
    ld $ra, 24($sp)
    daddiu $sp, $sp, 32
pair_return:
    jrc $ra
\end{minted}

\texttt{sd} stores all eight bytes of a register; \texttt{ld} loads all eight.
Saving with \texttt{sw} instead would retain only the low word. A later word load
could not recover a saved register's arbitrary upper half. These slots are naturally aligned: the frame base is a multiple of sixteen,
and the offsets are multiples of eight. This layout follows our calling
convention.

After saving the old values, \texttt{\$s0} holds the second input. The first call
returns its result in \texttt{\$v0}; \texttt{move \$s1, \$v0} keeps it across
the second call. That second call receives the retained input through
\texttt{\$a0}. On return, its result is copied into \texttt{\$v1} before the
first result is put back into \texttt{\$v0}. The two results are then independent
of the saved registers we restore.

Each inner \texttt{balc} replaces \texttt{\$ra}. The original continuation is
safe in its memory slot. Loading it back makes the final \texttt{jrc} return to
\texttt{main}, rather than to the continuation of an inner call. Restoring the
pointer last keeps all slot offsets relative to the same frame base.

Releasing a frame does not erase its bytes. It ends this function's use of those
slots. A later call may overwrite them, so an old saved value visible below the
current pointer is not evidence that the old frame is still active.

\section{Connect the functions to a complete program}

Save the following driver followed by \texttt{reverse\_pair} and
\texttt{reverse\_byte} in one file. The complete \texttt{reverse-rows.asm}
download is in \texttt{revision/labs/files}; it has exactly those three parts.
Select MIPS64 Release 6 and the Default memory layout.

\MARSCodeLabel{Prepare two inputs, call the pair routine and store its results}
\begin{minted}{mars-idea}
.data
.align 2
input_rows: .word 0x16, 0x81
results:    .word 0, 0
rows_guard: .word 0x13579bdf

.text
.globl main
main:
    dli $t0, -16
    and $sp, $sp, $t0
    la $s0, results
    la $t0, input_rows
    lw $a0, 0($t0)
    lw $a1, 4($t0)
main_call:
    balc reverse_pair
after_pair:
    sw $v0, 0($s0)
    sw $v1, 4($s0)
done:
    li $v0, 10
    syscall
\end{minted}

The Default layout initializes \texttt{\$sp} to \texttt{0x7FFFEFFC}. This is not
a 16-byte boundary. The driver loads the 64-bit mask $-16$, whose low four bits
are zero and whose other bits are one, and clears those pointer bits with
\texttt{and}. The resulting pointer is \texttt{0x7FFFEFF0}. This one-time setup
happens before our first call; it is not part of a function's per-call frame.
Every function still returns with the pointer it received.

Here \texttt{\$s0} is useful caller state: it points to the two result words.
The pair routine uses the same register for a different purpose, then restores
that pointer. The stores after the call can therefore use the caller's original
address. The second store has byte displacement four because each result object
is a word. Neither store should reach \texttt{rows\_guard}.

The program prints nothing. After a normal Run, inspect \texttt{results}: its
words are \texttt{0x00000068} and \texttt{0x00000081}. The input words and guard
are unchanged. The exit service reuses \texttt{\$v0}, so its final value is not
the function's returned first result. Inspect the return boundary or the stored
word when checking that result.

\section{Trace the lifetimes, then change the implementation}

Break at \texttt{main\_call} and record the result pointer and aligned stack
pointer. Use \emph{Step Into} to enter the call. At the function's first
instruction, \texttt{\$ra} names \texttt{after\_pair}; no frame has been reserved
yet. At \texttt{pair\_frame\_ready}, the pointer is 32 bytes lower and its three
slots contain the caller's saved values. The result pointer has not yet been
replaced by the next \texttt{move}.

At \texttt{pair\_first\_ready}, the first reversal is in \texttt{\$v0} and the
second input is in \texttt{\$s0}. After saving that first result and preparing
the second argument, step into the next call. Predict which values must remain
unchanged while the leaf routine works. At \texttt{pair\_return}, the results
are ready and both saved registers and \texttt{\$sp} have their incoming values.
The pending return should next reach \texttt{after\_pair}.

Try a deliberately faulty copy that omits the restore of \texttt{\$s0}. The
arithmetic can still be right, but the first result store now uses a row value as
an address. Another faulty copy can leave the first result only in
\texttt{\$v0}; the second call then replaces it. Identify the first lost value
rather than changing the bit-reversal loop in response to either symptom.

Finally, implement a separate \texttt{reverse\_and\_invert} routine. It receives
the same low-byte argument, uses \texttt{reverse\_byte}, and returns the reversed
byte with every one of its eight bits inverted. For \texttt{0x16}, reversing
gives \texttt{0x68}; XOR with \texttt{0xFF} gives \texttt{0x97}, decimal 151.

XOR produces one where two corresponding bits differ and zero where they agree.
With mask \texttt{0xFF}, it toggles each low-byte bit and leaves all higher bits
unchanged. The instruction \texttt{xori} supplies that mask as an immediate
operand; its destination and source come first, as with \texttt{andi}.

Keep the same preservation convention. Plan the value that must survive the call,
write the routine, then test all 256 byte inputs. The lab provides a caller and
public checks. The checks compare behavior and preservation; they do not require
your temporary registers or stack offsets to match this chapter.
